%-----------------------------------------------------------------------
\chapter{Wöhlerversuch}%\label{chap:kap1}
%-----------------------------------------------------------------------

Da für die Entwicklung und Validierung der MATLAB-Auswertesoftware nur eine begrenzte Anzahl realer Versuchsdaten zur Verfügung stand, wurden auf Grundlage der gemessenen Proben zusätzlich fiktive Versuchsdaten erstellt. Diese orientieren sich hinsichtlich ihrer Kraftniveaus, Schwingspielzahlen und Messsignalverläufe an den realen Versuchsergebnissen und dienen der Überprüfung der Funktionsfähigkeit sowie der Validierung der implementierten Auswertesoftware.

Zu Beginn des Kapitels wird eine kurze Bedienungsanleitung mit einer ergänzenden Checkliste vorgestellt, welche die erforderlichen Vorbereitungen für eine erfolgreiche Versuchsdurch- und eine erfolgreiche Programmausführung übersichtlich zusammenfasst und dem Nutzer als Orientierung dient.
	
	\section{Bedienungsanleitung für den Nutzer}
Um die Versuchsdurchführung normgerecht zu gestalten und die Auswertesoftware korrekt sowie gewinnbringend einzusetzen, muss der Versuch zuvor geplant, durchgeführt und dokumentiert werden. Die nachfolgende Bedienungsanleitung beschreibt stichpunktartig die Schritte vor, während und nach dem Versuch (bis zum Start der Auswertesoftware).

\textbf{Vorbereitung der Versuchsdurchführung:}
	\begin{itemize}
    \item Falls die Probenanzahl nicht extern vorgegeben ist, richtet sich diese nach der gewünschten Schätzgenauigkeit der Parameter der Grundgesamtheit für die Zeitfestigkeit ($N_{50\,\%,\mathrm{GG}}$, $s_{\lg N,\mathrm{GG}}$ und $k_{\mathrm{GG}}$) und für die Langzeitfestigkeit ($L_{aL,NG,\mathrm{GG}}$ sowie $s_{\lg L,\mathrm{GG}}$); für die zugehörigen Tabellen siehe Anhang~\ref{anhang:tabelle_erforderliche_Probenanzahl}.
    \item Vorgabe der Prüffrequenz und des Spannungsverhältnisses (deren Angabe erfolgt im Eingabefenster)~\cite[Seite~65]{DIN50100}
    \item Festlegung der Grenzschwingspielzahl anhand des Werkstoffs: für kubisch-raumzentrierte $N_{\text{G}} = 5 \times 10^6 \text{ Schwingspiele}$ und für kubisch-flächenzentrierte/hexagonal dichtest gepackte Werkstoffe $N_{\text{G}} = 10^7 \text{ Schwingspiele}$~\cite[Seite 16]{DIN50100}
    \item Bestimmung der geschätzten Langzeitfestigkeit (hierfür macht die DIN~50100 keine spezifischen Angaben, weshalb diese sinnvoll vom Benutzer gewählt werden muss; im Idealfall stützt sich die Schätzung auf Erfahrungswerte oder Vorversuche)~\cite[Seite~49]{DIN50100}
    \item Berechnen (mit \Eqref{eq:d_lg} und Anhang C der DIN~50100) oder Herauslesen des logarithmischen Stufensprungs (aus \Tabref{tab:Standardabweichung_Treppenstufenverfahren}) 
    \item Berechnung der Laststufen des Treppenstufenverfahrens (\Eqref{eq:Laststufen}): Bei Betrachtung der Tabelle in Anhang D.1.3 in der DIN 50100 liegt nahe, dass insgesamt zehn Laststufen ermittelt werden. Dabei werden in \Eqref{eq:Laststufen} für die Laufvariable beispielsweise $i=-4,\dots,0,\dots,5$ eingesetzt und die Resultate in der Tabelle im Anhang D.1.3 aus der DIN~50100 festgehalten.
    \item Allgemein empfiehlt es sich, die in der genannten Tabelle aufgeführten Parameter für die Versuchsplanung sowie die gesamte folgende Versuchsdurchführung zu dokumentieren~\cite[Seite~64--66]{DIN50100}. 
	\end{itemize}

\textbf{Während der Versuchsdurchführung:}
	  \begin{itemize}
	  \item Unmittelbare Vergabe der Probennummern in aufsteigender Reihenfolge für die Sicherstellung der Versuchsreihenfolge; dies ist für die Auswertung des Treppenstufenverfahrens zwingend erforderlich.
   	   \item Benennung der korrespondierenden Excel-Dateien muss folgendes Muster enthalten: z.B. $P\_{01}\_$ für Probe Nummer 1, $P\_{02}\_$ für Probe Nummer 2, etc.; \\ \textbf{Anmerkung}: Die zuerst durchgeführte Probe muss lediglich die kleinste Probennummer besitzen, kann jedoch auch bei einer anderen Probennummer ungleich eins starten.  
    	   \item Separierung der Daten in Zeitfestigkeitsbereich und Langzeitfestigkeitsbereich 
    	    \item Zudem ist folgende \textbf{Checkliste} während der Versuchsdurchführung zu beachten:
 \begin{itemize}[label=--]
     \item  Der Quotient der Schwingspielzahlen des höchsten und des niedrigsten gewählten Lasthorizontes beim Perlenschnurverfahren muss mindestens fünf betragen. Ein von der Norm anzustrebender Richtwert ist ein Quotient von 50. Unterschreitet dieser die Mindestvorgabe, so sind die Schätzungen von Mittelwert, Standardabweichung und Steigung nach DIN 50100 nicht aussagekräftig~\cite[Seite~38]{DIN50100}.
    \item Beim Treppenstufenverfahren müssen mindestens ein, aber maximal drei Lasthorizonte mit wenigstens einem Bruch und einem Durchläufer existieren. Weist ein Lasthorizont nur Brüche und ein zweiter nur Durchläufer auf, kann keine Standaradabweichung bestimmt werden. In diesem Fall ist gemäß \Eqref{eq:d_lg_1_2} ein neuer Stufensprung zu berechnen~\cite[Seite~48]{DIN50100}.
    \item Die Treppenstufenfolge erfordert mindestens zwei Umkehrpunkte und darf nicht unterbrochen werden~\cite[Seite~59]{DIN50100}
    \item Proben auf Laststufen, die anfänglich durchgeführt wurden und deren Treppenstufe im Laufe des Versuches nicht erneut erreicht werden, werden nicht gewertet~\cite[Seite~59]{DIN50100}; diese müssen nicht manuell entfernt werden sondern werden in der Auswertesoftware automatisch erkannt und nicht berücksichtigt.
    \item Für die entwickelte Auswertesoftware werden mindestens vier Proben im Zeitfestigkeitsbereich und zehn Proben im Langzeitfestigkeitsbereich vorausgesetzt. Diese Mindestanzahlen orientieren sich an den in der DIN 50100 angegebenen Tabellen zur statistischen Auswertung sowie an den in dieser Arbeit gewählten Auswerteverfahren (vgl. DIN 50100, Tabelle~2 bis Tabelle~4 sowie Tabelle~9 und Tabelle~10)~\cite[Seite~30--35 und Seite~68--69]{DIN50100}.
\end{itemize}
	\end{itemize}


	\textbf{Vorbereitung der Programmausführung:}
	   \begin{itemize}
	   \item Entfernen fehlerhafter Proben (z.B. beschädigter Excel-Dateien, vorzeitig abgebrochener Versuche, falsches Ausfallkriterium,...)
   	   \item Ausgangslaststufe $L_0$ festlegen
	\end{itemize}

		
	\section{Programmstruktur}
	\subsection{Programmaufbau und Funktionsstruktur}
	Die Auswertesoftware ist in ein Hauptprogramm und mehrere Funktionsdateien gegliedert. Die Aufgabe des Hauptprogramms besteht darin, die Funktionsdateien nacheinander aufzurufen. Dabei werden die erforderlichen Eingabedaten an die Funktionen übergeben und diese liefern daraufhin die gewünschten Rückgabeparameter.
	 Jede Funktion erfüllt eine klar definierte Aufgabe. Die verwendeten Funktionen und eine kurze Beschreibung ihrer Aufgaben sind in \Tabref{tab:funktionsuebersicht} zusammengefasst.
	
\begin{table}[H]
    \centering
    \caption{Übersicht der im Rahmen der Arbeit implementierten MATLAB-Funktionen}
    \label{tab:funktionsuebersicht}

    \begin{tabular}{|p{5.7cm}|p{8.5cm}|}
        \hline
        \textbf{Funktionsname} & \textbf{Aufgabe} \\
        \hline

        \texttt{Dateien\_einlesen} &
        Einlesen der Excel-Dateien für den Zeit- und Langzeitfestigkeitsbereich. \\
        \hline

        \texttt{Daten\_auslesen} &
        Extraktion der für die Auswertung relevanten Messgrößen aus den eingelesenen Messdateien. \\
        \hline
        
         \texttt{Abbruchkriterium} &
        Bestimmung der Bruchschwingspielzahl anhand des definierten Abbruchkriteriums. \\
        \hline

        \texttt{Kraftamplitude\_auslesen} &
        Ermittlung der Kraftamplitude aus den Messdaten. \\
        \hline


    \end{tabular}
\end{table}

Die Vorteile der Funktionsstruktur begründen die Entscheidung für diese strukturierte Arbeitsweise. Zum einen erhöhen Funktionen die Übersichtlichkeit, da durch die Auslagerung der Programmschritte das Hauptprogramm nicht überladen wird. Zum anderen ist der Code bei der Erstellung einfacher zu warten, falls ein Fehler auftritt. Des Weiteren werden die Erweiterbarkeit und Wiederverwendbarkeit des Programmcodes für den Laboringenieur erleichtert.

\clearpage

	\subsection{Programmablauf}
	\Figref{fig:programmablauf} zeigt den Ablauf des entwickelten MATLAB-Auswerteprogramms.
	
\begin{figure}[H]
    \centering

    \resizebox{\textwidth}{!}{%
    \begin{tikzpicture}[
        font=\fontsize{20}{24}\selectfont,
        node distance=13mm and 8mm,
        startstop/.style={
            ellipse,
            draw=black,
            thick,
            minimum width=4.6cm,
            minimum height=1.45cm,
            inner xsep=5mm,
            inner ysep=3mm,
            align=center
        },
        inputoutput/.style={
            trapezium,
            trapezium left angle=70,
            trapezium right angle=110,
            draw=black,
            thick,
            minimum width=6.2cm,
            minimum height=1.75cm,
            inner xsep=6mm,
            inner ysep=4mm,
            align=center
        },
        process/.style={
            rectangle,
            draw=black,
            thick,
            minimum width=5.8cm,
            minimum height=1.75cm,
            inner xsep=5mm,
            inner ysep=4mm,
            align=center
        },
        processwide/.style={
            rectangle,
            draw=black,
            thick,
            minimum width=7.2cm,
            minimum height=1.75cm,
            inner xsep=5mm,
            inner ysep=4mm,
            align=center
        },
        arrow/.style={
            thick,
            -{Latex[length=3.5mm]}
        }
    ]

    % Programmstart
    \node (start) [startstop]
        {Programmstart};

    \node (parameter) [inputoutput, below=of start]
        {Eingabe der Versuchs- und\\Auswerteparameter};

    \node (zfdateien) [inputoutput, below=of parameter]
        {Auswahl und Einlesen der\\Zeitfestigkeitsdaten};

    \node (lfdateien) [inputoutput, below=of zfdateien]
        {Auswahl und Einlesen der\\Langzeitfestigkeitsdaten};

    \node (aufteilung) [processwide, below=of lfdateien]
        {Getrennte Verarbeitung der Zeit- und\\Langzeitfestigkeitsdaten};

    % Zeitfestigkeitsbereich
    \node (zfaufbereitung) [
        process,
        below left=of aufteilung,
        xshift=-3mm
    ]
        {Auslesen und Aufbereiten der\\Zeitfestigkeitsdaten};

    \node (zfabbruch) [process, below=of zfaufbereitung]
        {Bestimmung des Abbruchkriteriums\\der Zeitfestigkeitsproben};

    \node (zfkraft) [process, below=of zfabbruch]
        {Ermittlung der Kraftamplituden\\der Zeitfestigkeitsproben};

    \node (zfauswertung) [process, below=of zfkraft]
        {Auswertung des\\Zeitfestigkeitsbereichs};

    % Langzeitfestigkeitsbereich
    \node (lfaufbereitung) [
        process,
        below right=of aufteilung,
        xshift=3mm
    ]
        {Auslesen und Aufbereiten der\\Langzeitfestigkeitsdaten};

    \node (lfabbruch) [process, below=of lfaufbereitung]
        {Bestimmung des Abbruchkriteriums\\der Langzeitfestigkeitsproben};

    \node (lfkraft) [process, below=of lfabbruch]
        {Ermittlung der Kraftamplituden\\der Langzeitfestigkeitsproben};

    \node (lfauswertung) [process, below=of lfkraft]
        {Auswertung des\\Langzeitfestigkeitsbereichs};

    % Zusammenführung
    \node (zusammenfuehrung) [
        processwide,
        below=12mm of $(zfauswertung)!0.5!(lfauswertung)$
    ]
        {Zusammenführung der Auswertungsergebnisse};

    \node (woehler) [processwide, below=of zusammenfuehrung]
        {Berechnung und Darstellung der Wöhlerlinie};

    \node (pdf) [processwide, below=of woehler]
        {Erstellung des PDF-Prüfprotokolls};

    \node (ende) [startstop, below=of pdf]
        {Programmende};

    % Gemeinsamer Ablauf
    \draw [arrow] (start) -- (parameter);
    \draw [arrow] (parameter) -- (zfdateien);
    \draw [arrow] (zfdateien) -- (lfdateien);
    \draw [arrow] (lfdateien) -- (aufteilung);

    % Verzweigung in die beiden Auswertungsbereiche
    \coordinate (verzweigung)
        at ($(aufteilung.south)+(0,-8mm)$);

    \draw [thick]
        (aufteilung.south) -- (verzweigung);

    \draw [arrow]
        (verzweigung) -| (zfaufbereitung.north);

    \draw [arrow]
        (verzweigung) -| (lfaufbereitung.north);

    % Zeitfestigkeitsbereich
    \draw [arrow] (zfaufbereitung) -- (zfabbruch);
    \draw [arrow] (zfabbruch) -- (zfkraft);
    \draw [arrow] (zfkraft) -- (zfauswertung);

    % Langzeitfestigkeitsbereich
    \draw [arrow] (lfaufbereitung) -- (lfabbruch);
    \draw [arrow] (lfabbruch) -- (lfkraft);
    \draw [arrow] (lfkraft) -- (lfauswertung);

    % Zusammenführung der beiden Bereiche
    \draw [arrow]
        (zfauswertung.south) |- (zusammenfuehrung.west);

    \draw [arrow]
        (lfauswertung.south) |- (zusammenfuehrung.east);

    % Abschließender Programmablauf
    \draw [arrow] (zusammenfuehrung) -- (woehler);
    \draw [arrow] (woehler) -- (pdf);
    \draw [arrow] (pdf) -- (ende);

    \end{tikzpicture}%
    }

    \caption{Schematischer Ablauf des entwickelten MATLAB-Auswerteprogramms}
    \label{fig:programmablauf}
\end{figure}


Mittels eines Eingabefeldes erfolgt zunächst eine Abfrage der Parameter für die Versuchsauswertung sowie für die Erstellung des Prüfberichts. Nach der Eingabe werden die entsprechenden Dateien der Versuchsproben für den Zeitfestigkeitsbereich und anschließend für den Langzeitfestigkeitsbereich in das Programm geladen. Das Programm arbeitet im weiteren Verlauf mit den getrennten Datensätzen der beiden Bereiche. Die benötigten Messdaten (Zeit, Weg und Kraft) sowie die Schwingspielzahl werden aus den Dateien extrahiert und gespeichert. Parallel wird anhand des Abbruchkriteriums die Bruchschwingspielzahl der Proben bestimmt. Daran schließt sich die Ermittlung der Kraftamplituden aus den Messdaten an. Anschließend erfolgt die Auswertung des Zeit- und Langzeitfestigkeitsbereichs. Aus den Ergebnissen der beiden Abschnitte wird daraufhin die Wöhlerlinie erstellt und in einem Diagramm visualisiert. Abschließend werden alle wesentlichen Auswertungsergebnisse in das Prüfprotokoll eingetragen und für den Benutzer bereitgestellt.   

Die einzelnen Schritte des Programmablaufs werden in den folgenden Kapiteln näher erläutert. 

	\section{Eingabe der Versuchsparameter}
	\subsection{Gestaltung des MATLAB-App-Eingabefensters}
	Für eine benutzerfreundliche Bedienung wurde das Eingabefenster mit dem MATLAB-App Designer erstellt. Die App erfasst alle für die Versuchsauswertung und die Erstellung des Prüfprotokolls erforderlichen Parameter.
	 Ein wesentlicher Vorteil der App besteht darin, dass das Eingabefenster jederzeit um Erfassungsfelder erweitert werden kann, ohne Änderungen am Hauptprogramm vornehmen zu müssen. Durch den gezielten Einsatz von Eingabefeldern, Dropdown-Menüs und vorgegebenen Standardwerten werden Fehleingaben reduziert, was die Dateneingabe erheblich vereinfacht und den zeitlichen Aufwand für den Anwender reduziert.
	 
	 Für die Eingabe werden Textfelder, numerische Eingabefelder, Dropdown-Menüs sowie Optionsfelder verwendet. Nach Fertigstellung der Eingaben werden über den  \glqq Auswahl bestätigen\grqq{}-Button die eingetragenen Werte an das Hauptprogramm übergeben. 
	 
	
	\subsection{Eingabeparameter}
	Das Eingabefenster ist entsprechend seiner Funktion in zwei Bereiche unterteilt.
	Im oberen Bereich befinden sich die Versuchsdaten, die in das PDF-Prüfprotokoll eingetragen werden. Dazu zählen der Prüfer, das Datum, der Kunde, die allgemeine Probenbezeichnung, die Prüftemperatur und die Frequenz der Schwingprüfung.
	Der untere Bereich enthält die Werkstoffwahl, die Grenzschwingspielzahl, den Wöhlerlinien-Typ, den Lasthorizont $L_0$ sowie die Standardabweichung der Grundgesamtheit. 
	
	Das Eingabefenster der entwickelten MATLAB-App ist in \Figref{fig:Eingabefenster} dargestellt.
	
	 \begin{figure}[H]
        \centering
            \includegraphics[width=1.0\textwidth]{Eingabefenster_MATLAB.png}
            \caption{Eingabefenster MATLAB}
  	  \label{fig:Eingabefenster}
\end{figure}
	
	Zusätzlich sorgt das Eingabefenster für automatische Vorschläge. Die Werkstoffauswahl dient hierbei als Grundlage der standardisierten Vorgabe der Grenzschwingspielzahl und des Wöhlerlinien-Typs gemäß den Empfehlungen der DIN 50100. Die vorgeschlagenen Werte können vom Benutzer übernommen oder bei Bedarf angepasst werden. Der Wöhlerlinien-Typ bestimmt die unterschiedliche Auswertung des Langzeitfestigkeitsbereichs und die Grenzschwingspielzahl legt fest, ab wann eine Probe als Durchläufer gewertet wird. Darüber hinaus werden für das Prüfdatum, die Prüftemperatur und die Prüffrequenz ebenfalls Standardwerte bereitgestellt. 
	
	Der Lasthorizont $L_0$ stellt die Ausgangslaststufe des Treppenstufenverfahrens dar. Die Standardabweichung dient der Berechnung des logarithmischen Stufensprungs zwischen den einzelnen Laststufen.  
	
	Über \texttt{getOutputData} werden die eingegebenen Werte an das Hauptprogramm übergeben und den nachfolgenden Programmfunktionen zur Verfügung gestellt. 

	\section{Einlesen der Messdaten}
	Das Einlesen der Messdaten beginnt mit der Dateiauswahl durch den Benutzer. Akzeptierte Dateien sind Excel-Dateien im Format \texttt{.xlsx} und \texttt{.xls}. Aufgrund der unterschiedlichen Auswerteverfahren erfolgt das Importieren der Messdaten der Zeitfestigkeit und der Langzeitfestigkeit getrennt. Der Nutzer wird folglich mit einem Informationsfenster dazu aufgefordert, zuerst die Dateien des Zeitfestigkeitsbereiches auszuwählen. Erst nach Abschluss dieser Auswahl erscheint ein weiteres Dialogfenster für den Langzeitfestigkeitsbereich, um die zugehörigen Dateien auszuwählen. 
 Das Informationsfenster wird durch die MATLAB-Funktion \texttt{msgbox} erzeugt. 
 
 Die eigentliche Dateiauswahl erfolgt mithilfe der MATLAB-Funktion \texttt{uigetfile}, mit der mehrere Dateien gleichzeitig in das Programm eingelesen werden können. Eine \texttt{while}-Schleife sorgt dafür, dass sich der Dateiauswahldialog so lange erneut öffnet, bis der Nutzer keine weitere Datei auswählt. Zur Vermeidung redundanten Programmcodes wurde der gesamte Dateieneinlesevorgang in die Funktion \texttt{Dateien\_einlesen} ausgelagert.
 
 Nach Abschluss der Dateiauswahl wird überprüft, ob für beide Auswertebereiche eine ausreichende Anzahl an Proben vorliegt. Für den Zeitfestigkeitsbereich müssen mindestens vier Proben und für den Langzeitfestigkeitsbereich mindestens zehn Proben ausgewählt worden sein. Wird diese Mindestanzahl unterschritten, gibt das Programm eine entsprechende Fehlermeldung aus bricht die weitere Programmausführung ab, da für kleinere Stichproben nach DIN~50100 keine Fehler zur Grundgesamtheit angegeben sind.
 
 Die mit \texttt{readcell} eingelesenen Tabellen werden jeweils als vollständiges Cell-Array in einem übergeordneten Cell-Array gespeichert. Jedes Element dieses übergeordneten Cell-Arrays enthält damit die vollständigen Rohmessdaten einer Probe in der ursprünglichen Tabellenstruktur der Hydropulsanlage. Dies erleichtert die spätere Weiterverarbeitung der einzelnen Proben. Die Dateien des Zeitfestigkeits- und des Langzeitfestigkeitsbereichs werden dabei in zwei getrennten Cell-Arrays gespeichert. Im nächsten Schritt werden alle für die Versuchsauswertung relevanten Messgrößen aus den Proben extrahiert und in neuen Variablen abgelegt.

	\section{Aufbereitung der Messdaten}
	
	\subsection{Extraktion der relevanten Messgrößen}
	
	Das Ziel dieses Schrittes ist es aus den Tabellen der einzelnen Proben alle relevanten Daten zu extrahieren. Dieser Verarbeitungsschritt wird durch die Funktion \texttt{Daten\_auslesen} ausgeführt. Für die Auswertung werden die Messgrößen Kraft, Weg, Zeit und Schwingspielzahl benötigt. Zur Identifikation dieser wird zunächst die Zeile, die das Wort Zeit enthält lokalisiert. Ausgehend von dieser Zeile werden anschließend die Spaltenüberschriften für Kraft, Weg und Schwingspielzahl ermittelt. Dadurch wird eine Suche unabhängig der Spaltenreihenfolge gewährleistet. 
	
	In \Tabref{tab:sensoren_uebersicht} sind die verschiedenen Kraft- und Wegsensoren der drei Hydropulsanlagen des Maschinendynamiklabors aufgelistet. Jede Hydropulsanlage verfügt über einen Kraft- als auch einen Wegsensor, der fest an der Maschine verbaut ist. Zusätzlich ist ein Lasersensor, der die Wegänderung erfasst, vorhanden und kann flexibel bei allen drei Anlagen eingesetzt werden.

\begin{table}[ht]
    \centering
    \caption{Übersicht der Sensoren und deren Zuordnung zu den Hydropulsanlagen}
    \label{tab:sensoren_uebersicht}
    \small
    \renewcommand{\arraystretch}{1.2} % Erhöht den Zeilenabstand für bessere Lesbarkeit
    \begin{tabular}{llll}
        \hline
        \textbf{Bezeichnung} & \textbf{Messgröße} & \textbf{Zuordnung Anlage} & \textbf{Status} \\
        \hline
        Weg PL100K   & Weg [mm]   & Hydropulsanlage 1 & Fest verbaut \\
        Kraft PL100K & Kraft [kN] & Hydropulsanlage 1 & Fest verbaut \\
        \hline
        Weg PL16K   & Weg [mm]   & Hydropulsanlage 2 & Fest verbaut \\
        Kraft PL16K & Kraft [kN] & Hydropulsanlage 2 & Fest verbaut \\
        \hline
        Weg PLF7D  & Weg [mm]   & Hydropulsanlage 3 & Fest verbaut \\
        Kraft PLF7D & Kraft [kN] & Hydropulsanlage 3 & Fest verbaut \\
        \hline
        Weg Laser   & Weg [mm]   & Flexibel (Anlage 1--3) & Wechselbar \\
        \hline
    \end{tabular}
\end{table}

Um den MATLAB-Code für zukünftige Versuche an allen drei Hydropulsnalagen zu standardisieren, werden alle Sensoren in den Code mit einbezogen. 

\Figref{fig:weg_laser} zeigt den Vergleich der aufgezeichneten Wegsignale des PL16K-Sensors und des Lasersensors.
 Wie aus \Figref{fig:weg_laser} hervorgeht, weist das Wegsignal des Sensors PL16K im Vergleich zum Lasersensor deutlich stärkere Oszillationen auf. Die Ursache hierfür liegt in der direkten Montage des PL16K-Sensors am Zylinder. Dadurch erfasst er mechanische Schwingungen und Vibrationen der Zylinderbewegung direkt im Messsignal. 
   Der Lasersensor hingegen, ist an einer separaten, unbewegten Stelle befestigt. Dadurch liefert dieser ein ruhigeres und weniger von den Schwingungen der Hydropulsanlage beeinflusstes Wegsignal. Aufgrund dieser Eigenschaft prüft das Programm zunächst, ob der Lasersensor verfügbar ist. Ist dies der Fall wird dessen Messsignal anschließend für die Auswertung verwendet. Andernfalls greift das Programm auf den Wegsensor der jeweiligen Hydropulsanlage zurück.  
	
	    \begin{figure}[H]
        \begin{center}
            \includegraphics[width=1.0\textwidth]{Vergleich der Wegsignale.pdf}
            \caption{Vergleich Wegsignal des Wegaufnehmers Lasersensor und des Wegaufnehmers PL16K.}
  	  \label{fig:weg_laser}
        \end{center}
    \end{figure}

	\subsection{Vereinheitlichung und Speicherung der Messdaten}
	Die extrahierten Messgrößen Kraft, Weg, Zeit und Schwingspielzahl werden in separate Matrizen gespeichert. Dabei entspricht jede Spalte den Messwerten einer Probe. Für eine einheitliche Weiterverarbeitung werden die in den Rohmessdaten in Kilonewton vorliegenden Kraftwerte in Newton umgerechnet. Da der \texttt{QuarterCycleCounter} stets den vierfachen Wert der Schwingspielzahl ausgibt, werden die ausgelesenen Werte durch vier dividiert. Die Daten des Zeit- und des Langzeitfestigkeitsbereiches werden zunächst getrennt gespeichert.  
	
	\section{Bestimmung der Bruchschwingspielzahl}
	Nach der Aufbereitung und Vereinheitlichung der Messdaten erfolgt im nächsten Schritt die Bestimmung der Bruchschwingspielzahl. Hierfür wird die Programmfunktion \texttt{Abbruchkriterium} verwendet.
	 
	 Da die Hydropulsanlage einer gewissen Regelverzögerung unterliegt und dadurch ein Nachschwingen erfolgen kann ist das Ziel des Abbruchkriteriums eine reproduzierbare Bruchschwingspielzahl zu bestimmen, die möglichst nicht durch das Nachschwingen der Hydropulsanlage beeinflusst wird. Dies dient vorallem der Vergleichbarkeit der einzelnen Proben.
	 Das Abbruchkriterium stützt sich auf die Wegänderung, da bei Probenversagen ein sprunghafter Anstieg des Weges erfolgt, wodurch sich der Bruch eindeutig erkennen lässt. Da die Rohmessdaten nur die letzten 30 Sekunden vor der Beendigung des Versuchs enthalten, weist das Wegsignal zu diesem Zeitpunkt bereits einen individuellen Regeloffset auf. Es kann daher kein absoluter Wegwert als allgemeines Abbruchkriterium verwendet werden. Des Weiteren unterscheidet sich das Bruchverhalten von Werkstoff zu Werkstoff aufgrund unterschiedlicher Werkstoff- und Bauteileigenschaften, weshalb eine Weggrenzwert als Abbruchkriterium riskant wäre. Aus diesem Grund wird die Steigungsänderung betrachtet. Hier wurde anhand der vorliegenden Messdaten empirisch ein Faktor für die Steigungsänderung von 5 bestimmt. Diese Steigungsänderung wird dabei auf die Anfangssteigung des aufgezeichneten Messbereichs bezogen. 
	 Da die Wegänderung die periodischen Schwingungen enthält, kann die Steigung nicht aus zwei nebeneinader liegenden Messpunkten bestimmt werden, das die Steigung sonst die Steigung der einzelnen Schwingspiele darstellt. Aus diesem Grund wird das geglaettete Wegsignal der einzelnen Proben mit der MATLAB-Funktion \texttt{movmean} in einer vorgegebene Fenstergröße ermittelt. Die Fenstergröße entspricht hierbei der Anzahl der Messpunkte einer Schwingung. Diese ergeben sich aus der Division der Dauer einer Schwingung (1/Prüffrequenz) und dem aus den in eingelesenen Excel-Tabellen stehenden Abtastintervall. 
	 
	Da bei den Durchläufern kein Bruch auftritt und somit kein Weganstieg erfolgt, wird der letzte Messpunkt (die Endschwingspielzahl) verwendet. 
	
	Da sich das aufgezeichnete Wegsignal bei unterschiedlichen Werkstoffen und Probengeometrien unterscheiden kann, ist bei zukünftigen Versuchsreihen der verwendete Steigungsfaktor zu überprüfen und gegebenenfalls anzupassen. 
	Dazu wird das Ergebnis der Abbruchwerte in einem Weg-Zeit-Plot dargestellt. Der in der Auswertesoftware ermittelte Abbruchpunkt wird im Diagramm markiert. Dadurch kann der Nutzer diesen Vorgang leichter nachvollziehen und die Abbruchpunkte auf Richtigkeit überprüfen.
	
	 \begin{figure}[H]
        \centering
            \includegraphics[width=1.0\textwidth]{Diagramm_Abbruchkriterium.png}
            \caption{Darstellung des Abbruchpunkts im Weg-Zeit-Diagramm der untersuchten Proben.}
  	  \label{fig:Abbruchpunkt}
    \end{figure}
	
	 \begin{figure}[H]
        \centering
            \includegraphics[width=1.0\textwidth]{Diagramm_Abbruchkriterium_Detailansicht.png}
            \caption{Deitailansicht zur Darstellung des Abbruchpunkts im Weg-Zeit-Diagramm der untersuchten Proben.}
  	  \label{fig:Abbruchpunkt_Detail}
    \end{figure}
	
	Die ermittelte Bruchschwingspielzahl dient anschließend zusammen mit der Kraftamplitude als Grundlage für die Berechnung der Wöhlerlinie.
	
\textbf{Eingangs- und Ausgangsgrößen?
Wahl der Fenstergröße?
Übergabe der Ergebnisse an das Hauptprogramm. noch was zu den Abbildungen schreiben!}

	\section{Berechnung der Kraftamplitude}
	Die Kraftamplitude stellt die letzte benötigte Kenngröße zur Auswertung des Zeitfestigkeits- und der Langzeitfestigkeitsbereichs dar. Diese wird mithilfe der Teilfunktion des Auswerteprogramms \texttt{Kraftamplitude\_auslesen} ermittelt. 
	
	Da die Probe aufgrund der zyklischen Belastung einen sinusförmige Prüfkraft erfährt, ergibt sich die Kraftamplitude analog zu \Eqref{eq:Spannungsamplitude}. Aus diesem Grund genügt die Bestimmung der Minima und Maxima. 
	
	Zur Bestimmung der lokalen Maxima des Kraftsignals wird die MATLAB-Funktion \texttt{findpeaks} verwendet. Diese werden dann in einem zeitlichen Mindestabstand auch MinPeakDistance genannt gesucht. Dieser verhindert die Mehrfacherkennung innerhalb einer Schwingung. Der gewählte Mindestabstand berechnet sich aus der Prüffrequenz und entspricht 70\,\% der Dauer einer Schwingungsperiode, um mehrere Maxima innerhalb derselben Schwingung zu vermeiden. Da \texttt{findpeaks} ausschließlich Maxima erkennt wird zur Bestimmung der Minima das Kraftsignal invertiert und anschließend mit \texttt{findpeaks} ausgewertet. Die ermittelten Minima werden anschließend wieder auf ihr ürsprüngliches Vorzeichen zurückgesetzt. Aufgrund dessen kann diese Vorgehensweise unabhängig des gewählten Spannungsverhältnis R (=Lastverhältnis) angewendet werden. 
	
	Anschließend wird der Mittelwert aller Minima und aller Maxima berechnet, da die einzelnen Peaks aufgrund des vorhandenen Messrauschens sowie Reglerschwankungen leicht schwanken können. Deshalb liefert der Mittelwert eine robustere Bestimmung der Kraftamplitude.
	Aus der halben Differnz der beiden Mittelwerte ergibt sich anschließend die Kraftamplitude. 
	
	Die berechneten Kraftamplituden werden als Vektor an das Hauptprogramm zurückgegeben. Dabei wird ein Kraftamplitudenwert pro Probe ausgegeben. 
	
	\textbf{warum 0,7*Dauer Schwingung? Begründung schreiben!}
	
	
	\clearpage
	
	\section{Auswertung des Zeitfestigkeitsbereichs}
	Mit Hilfe der ermittelten Kraft- bzw. Spannungsamplituden und der Bruchschwingspielzahlen der Zeitfestigkeitsproben kann der Zeitfestigkeitsbereich ausgewertet werden. Hierzu wird das in \Chapref{chap:Theorie_Zeitfestigkeit} vorgestellte Perlenschnurverfahren im MATLAB-Programm realisiert. Abhängig davon, ob der belastete Querschnitt angegeben wurde, erfolgt die Auswertung anhand der Kraft- oder Spannungsamplitude. Diese werden im Folgenden (analog zur korrespondierenden Variable im MATLAB-Programm) zusammenfassend als Auswerteamplitude bezeichnet.
	Ziel dieses Auswerteschrittes ist die Bestimmung der 50\,\%-Zeitfestigkeitsgeraden sowie ihrer Parameter $C$ und $k$. Diese werden im weiteren Programmablauf gemeinsam mit der Langzeitfestigkeit zur Ermittlung der Knickschwingspielzahl verwendet.
	Die nachfolgend beschriebenen Auswertungsschritte sind im MATLAB-Programm implementiert.
	
	
	\subsection{Berechnung der Regressionsgeraden}
	Zur Veranschaulichung werden die einzelnen Rechenschritte im weiteren Verlauf dieses Abschnitts exemplarisch anhand der generierten Proben demonstriert. Als Ausgangsgrößen dienen die zuvor vom MATLAB-Programm ermittelten Kraftamplituden und Bruchschwingspielzahlen der Zeitfestigkeitsproben, die in \Tabref{tab:Zeitfestigkeitsproben} aufgeführt sind.
	\begin{table}[H]
    \centering
    \caption{Kraftamplituden und Bruchschwingspielzahlen der Zeitfestigkeitsproben}
    \label{tab:Zeitfestigkeitsproben}
    \begin{tabular}{c|c|c}
        \hline
        \textbf{Probe} & 
        \textbf{Kraftamplitude $F_a$ in N} & 
        \textbf{Bruchschwingspielzahl $N_B$} \\
        \hline
        1  & 2313,2 & 71\,098 \\
        2  & 1599,8 & 1\,533\,291 \\
        3  & 2195,7 & 101\,820 \\
        4  & 2075,9 & 144\,990 \\
        5  & 1926,2 & 219\,990 \\
        6  & 1776,5 & 334\,990 \\
        7  & 1656,8 & 469\,993 \\
        8  & 1547,0 & 649\,990 \\
        9  & 1447,2 & 892\,550 \\
        10 & 1417,3 & 979\,990 \\
        \hline
    \end{tabular}
\end{table}
	
	Für die beispielhafte, ausführliche Rechnung wird lediglich Probe 1 mit der Kraftamplitude $F_a = 2313{,}2\,\mathrm{N}$ und Bruchschwingspielzahl $N_B = 71\,098$ herangezogen.
	
	Da die Basquin-Gleichung nach \Eqref{eq:Zeitfestigkeitsgerade_C} durch den Zehner-Logarithmus linearisiert wird, können die Parameter der Zeitfestigkeitsgeraden mittels linearer Regression bestimmt werden. Hierfür werden die Auswerteamplitude und die Bruchschwingspielzahl gemäß \Eqref{eq:substitution} logarithmiert. Im Fall von Probe 1 resultiert:
\begin{align*}
x_1 &= \lg(F_{a,1}) = 3{,}3642, \\
y_1 &= \lg(N_{B,1}) = 4{,}8519.
\end{align*}
	
	 Im MATLAB-Code werden die logarithmierten Werte aller Proben analog in den Variablen x und y gespeichert. Aus diesen wird zudem das Produkt $xy$ sowie das Quadrat aus $x$ gebildet:
\begin{align*}	 
x_1 y_1 &= 3{,}3642 \cdot 4{,}8519 = 16{,}3248, \\
x_1^2 &= (3{,}3642)^2 = 11{,}3179.
\end{align*}

Führt man dies entsprechend für alle Proben des Zeitfestigkeitsbereiches durch, ergibt sich folgende \Tabref{tab:Regressionswerte_Zeitfestigkeit}:
\begin{table}[H]
    \centering
    \caption{Logarithmierte Versuchswerte und Hilfsgrößen der Regressionsrechnung}
    \label{tab:Regressionswerte_Zeitfestigkeit}
    \begin{tabular}{c|c|c|c|c}
        \hline
        \textbf{Probe} &
        $\boldsymbol{x = \lg F_a}$ &
        $\boldsymbol{y = \lg N_B}$ &
        $\boldsymbol{xy}$ &
        $\boldsymbol{x^2}$ \\
        \hline
        1  & 3,3642 & 4,8519 & 16,3227 & 11,3180 \\
        2  & 3,2041 & 6,1856 & 19,8191 & 10,2660 \\
        3  & 3,3416 & 5,0078 & 16,7341 & 11,1661 \\
        4  & 3,3172 & 5,1613 & 17,1213 & 11,0039 \\
        5  & 3,2847 & 5,3424 & 17,5482 & 10,7893 \\
        6  & 3,2496 & 5,5250 & 17,9540 & 10,5597 \\
        7  & 3,2193 & 5,6721 & 18,2600 & 10,3637 \\
        8  & 3,1895 & 5,8129 & 18,5402 & 10,1728 \\
        9  & 3,1605 & 5,9506 & 18,8071 & 9,9889 \\
        10 & 3,1514 & 5,9912 & 18,8810 & 9,9316 \\
        \hline
    \end{tabular}
\end{table} 
	 
	 Im Anschluss werden die für die Regressionsrechnung erforderlichen logarithmierten Werte aufsummiert:
\begin{align*}
\sum x   &= 3{,}3642 + 3{,}2041 + 3{,}3416 + \ldots = 32{,}4821, \\
\sum y   &= 4{,}8519 + 6{,}1856 + 5{,}0078 + \ldots = 55{,}5008, \\
\sum xy  &= 16{,}3228 + 19{,}8193 + 16{,}7341 + \ldots = 179{,}9877, \\
\sum x^2 &= 11{,}3178 + 10{,}2663 + 11{,}1663 + \ldots = 105{,}5600.
\end{align*}

Mit diesen Summen und der Probenanzahl $n$ werden die Regressionsparameter $a_0$ und $a_1$ nach \Eqref{eq:reg_a1} und \Eqref{eq:reg_a0} berechnet: 
\begin{align*}
a_1
&= \frac{n\sum xy-\sum x\sum y}
        {n\sum x^2-\left(\sum x\right)^2} \\[1mm]
&= \frac{10\cdot179{,}9877-32{,}4821\cdot55{,}5008}
        {10\cdot105{,}5600-(32{,}4821)^2} \\
&= -5{,}6618,
\end{align*}

\begin{align*}
a_0
&= \frac{\sum y-a_1\sum x}{n} \\[1mm]
&= \frac{55{,}5008-(-5{,}6618)\cdot32{,}4821}{10} \\
&= 23{,}7709.
\end{align*}

In MATLAB erhält man für $\sum x$, $\sum xy$ und $\sum x^2$ identische Werte, während für $\sum y = 55.5010$, $a_0 = 23.8768$ und $a_1 = -5.6421$ minimale Abweichungen auftreten. Diese sind auf die Rundung der dargestellten Werte zurückzuführen, da MATLAB intern mit einer höheren numerischen Genauigkeit rechnet. Um präzisere Ergebnisse für die Kenngrößen $k$ und $C$ zu erreichen, werden nachfolgend die MATLAB-Resultate verwendet. 

	Aus $a_0$ und $a_1$ werden daraufhin der Steigungsparameter $k$ und der Lageparameter $C$ der Zeitfestigkeitsgeraden nach \Eqref{eq:reg_C} und \Eqref{eq:reg_k} bestimmt: 
\begin{align*}
k &= -a_1 = -(-5{,}6421) = 5{,}6421, \\
C &= 10^{a_0} = 10^{23{,}8768} = 7{,}53 \cdot 10^{23}.
\end{align*}
	Die durch $C$ und $k$ beschriebene Regressionsgerade entspricht der 50\,\%-Zeitfestigkeitsgeraden. Für deren Darstellung werden die zugehörigen Schwingspielzahlen an der höchsten und der niedrigsten Auswerteamplitude ermittelt. 
	
	\subsection{Berechnung der Standardabweichung und der Streulinien}
	Zur Bestimmung der 10\,\%- und 90\,\%-Ausfallwahrscheinlichkeitslinien muss zunächst die Streuung der Versuchsergebnisse anhand der Standardabweichung ermittelt werden.
	Gemäß DIN~50100 werden die Versuchsergebnisse zur Bestimmung der Standardabweichung auf einen fiktiven Lasthorizont verschoben. Da dessen Lage frei wählbar ist, wird dieser im MATLAB-Programm als arithmetischer Mittelwert aus der höchsten und der niedrigsten Auswerteamplitude festgelegt.
Die fiktive Schwingspielzahl, die sich durch die parallele Verschiebung auf den fiktiven Lasthorizont ergibt, wird nach \Eqref{eq:N_fiktiv} berechnet.  
	
Zur Veranschaulichung der Berechnung der Standardabweichung ist in \Figref{fig:verschieben_Lasthorizonte} dargestellt, wie die einzelnen Proben parallel zur Zeitfestigkeitsgeraden auf den fiktiven Lasthorizont verschoben werden.
	
	\begin{figure}[H]
        \centering
            \includegraphics[width=1.0\textwidth]{verschieben_Lasthorizonte.png}
            \caption{Verschieben der Versuchsergebnisse auf einen gemeinsamen Lasthorizont (In Anlehnung an DIN~50100, Seite 		56).}
  	  \label{fig:verschieben_Lasthorizonte}
    \end{figure}

Aus den logarithmierten, fiktiven Schwingspielzahlen wird als nächstes gemäß \Eqref{eq:mittelwert} der Mittelwert auf dem fiktiven Lasthorizont bestimmt. Darauf aufbauend wird die Standardabweichung nach \Eqref{eq:S_Perlenschnurverfahren} berechnet und anschließend mit dem in \Eqref{eq:s_korrektur} angegebenen Korrekturfaktor angepasst.

Mit der korrigierten Standardabweichung werden anschließend die Schwingspielzahlen für eine Ausfallwahrscheinlichkeit von 10\,\% und 90\,\% auf dem fiktiven Lasthorizont nach \Eqref{eq:N10} und \Eqref{eq:N90} bestimmt. Die 50\,\%-Zeitfestigkeitsgerade wird daraufhin parallel durch die entsprechenden Punkte verschoben, wodurch die 10\,\%- und 90\,\%-Zeitfestigkeitsgeraden entstehen.

Dienen weniger als 10 Proben zur Auswertung oder ist der Quotient der Schwingspielzahlen der 50\,\%-Zeitfestigkeitsgeraden am niedrigsten und höchsten untersuchten Lasthorizont kleiner als fünf, gibt das Programm eine entsprechende Warnmeldung aus.. 
Diese geben dem Nutzer die Information aus, dass die Schätzung des Mittelwertes, der Standardabweichung und der Steigung der Zeitfestigkeitsgerade statistisch nicht mehr erwartungsgetreu und unzuverlässig ist.

Die damit ermittelten Parameter C, k und die drei Zeitfestigkeitsgeraden werden anschließend mit den Ergebnissen des Langzeitfestigkeitsbereichs zusammengeführt.

\clearpage

	\section{Auswertung des Langzeitfestigkeitsbereichs}
	Die Auswertung des Langzeitfestigkeitsbereichs erfolgt nach dem in \Chapref{chap:Theorie_Treppenstufenverfahren} beschriebenen Treppenstufenverfahren. Ziel ist die Bestimmung der Langzeitfestigkeit einschließlich ihrer Streuung, die anschließend gemeinsam mit den Ergebnissen des Zeitfestigkeitsbereichs zur Ermittlung der Knickschwingspielzahl verwendet wird.
	
	\subsection{Zuordnung der Versuchsergebnisse zu den Laststufen}
	Für die statistische Auswertung des Treppenstufenverfahrens müssen die Versuchsergebnisse zunächst den entsprechenden Laststufen zugeordnet werden.
	Dazu werden die Brüche und Durchläufer auf den einzelnen Lasthorizonten zusammengezählt und als Ereignisse definiert. 
	
	Zunächst werden die Langzeitfestigkeitsproben anhand der ermittelten Schwingspielzahl in Brüche und Durchläufer eingeteilt.
	Hierzu wird die Schwingspielzahl jeder Probe mit der vorgegebenen Grenzschwingspielzahl verglichen. Liegt sie unterhalb der Grenzschwingspielzahl, wird die Probe als Bruch gewertet. Bei einer Schwingspielzahl größer oder gleich der Grenzschwingspielzahl handelt es sich um einen Durchläufer. 
	
	Wurde eine belastete Querschnittsfläche angegeben, wird der vorgegebene Lasthorizont $L_0$ analog zu den Auswerteamplituden von einer Kraft- in eine Spannungsamplitude umgerechnet. So können der Lasthorizont $L_0$ und die Auswerteamplituden direkt miteinander verrechnet werden.
	
	Da die tatsächlich gemessene Auswerteamplitude geringfügig vom vorgegebenen Lasthorizont $L_0$ abweichen kann, wird zunächst die gemessene Auswerteamplitude mit der geringsten betragsmäßigen Abweichung zu $L_0$ bestimmt. Diese wird im weiteren Programmablauf als gemessener Ausgangslasthorizont $L_{0,\mathrm{gemessen}}$ verwendet.
	
	Wurde für die Versuchsdurchführung der Wurzelstufensprung aus \Eqref{eq:d_lg_1_2} verwendet, wird die logarithmische Schrittweite für die Laststufenzuordnung entsprechend angepasst. Hierzu wird der eingegebene Wert der Standardabweichung halbiert. 
	
	Zur Zuordnung der Versuchsergebnisse wird \Eqref{eq:Laststufen} nach der Laststufennummer i umgestellt. Da die gemessenen Auswerteamplituden aufgrund von Messabweichungen nicht exakt den theoretischen Laststufen entsprechen, wird das berechnete Ergebnis auf die nächstgelegene ganzzahlige Laststufennummer gerundet.


	Anschließend werden nicht auswertbare Lasthorizonte entfernt. Zum einen werden die Proben, die eine Laststufennummer $i < 0$ aufweisen von der weiteren Auswertung ausgeschlossen, da sie unter dem untersten, gewerteten Lasthorizont $L_0$ liegen. Zum anderen werden nicht zu wertende Anfangsversuche nicht berücksichtigt. Gemäß der in \Chapref{chap:Theorie_Treppenstufenverfahren} beschriebenen Vorgehensweise werden zu Beginn der Treppenstufenfolge nur diejenigen Laststufen berücksichtigt, die im weiteren Versuchsverlauf erneut erreicht wurden. Hierzu überprüft das MATLAB-Programm die Laststufen vom höchsten Niveau abwärts und ermittelt, ab welcher Laststufe mehr als ein Ereignis vorliegt.
	
	Anschließend ergänzt das MATLAB-Programm gemäß der in \Chapref{chap:Theorie_Zeitfestigkeit} beschriebenen Vorgehensweise nach dem letzten realen Versuch eine fiktive Folgeprobe. Die zugehörige Laststufe wird dabei in Abhängigkeit vom Ergebnis der letzten realen Probe automatisch festgelegt.
	
	Anschließend überprüft das Programm, ob die ermittelte Laststufenfolge den Vorgaben des Treppenstufenverfahrens entspricht. Nach einem Durchläufer muss die folgende Probe auf der nächsthöheren und nach einem Bruch auf der nächstniedrigeren Laststufe liegen. Die daraus erwartete Laststufe wird mit der aus der gemessenen Auswerteamplitude bestimmten Laststufe verglichen. Bei einer Abweichung gibt das Programm eine Fehlermeldung aus
	 
	 Ist die Prüfung positiv wird anschließend die Treppenstufenfolge zur Veranschaulichung geplottet. 
	
    \begin{figure}[H]
        \centering
            \includegraphics[width=0.8\textwidth]{Treppenstufenfolge.png}
            \caption{Folge im Treppenstufenverfahren.}
  	  \label{fig:Treppenstufenfolge}
    \end{figure}

	\section{Automatisierte Erstellung des Prüfprotokolls}
Nach Abschluss der Versuchsauswertung werden die wesentlichen Prüfparameter und Ergebnisse automatisiert in einem Prüfprotokoll zusammengeführt. Hierzu verwendet das MATLAB-Programm eine zuvor erstellte LaTeX-Vorlage, deren Platzhalter während der Programmausführung durch die eingegebenen und berechneten Werte ersetzt werden.

Zunächst wird der Speicherort der LaTeX-Vorlage ermittelt und deren Inhalt als Text in MATLAB eingelesen. Kann die Vorlage nicht gefunden werden, wird die weitere Erstellung des Prüfprotokolls mit einer Fehlermeldung abgebrochen.

Die zuvor erstellte Wöhlerkurve wird mit exportgraphics als PDF-Datei im Vektorformat exportiert. Durch die vektorbasierte Ausgabe bleibt die Darstellung auch bei einer Skalierung im Prüfprotokoll ohne Qualitätsverlust erhalten.

Die LaTeX-Vorlage enthält definierte Platzhalter für allgemeine Angaben, Prüfparameter und Versuchsergebnisse. Diese werden mithilfe der MATLAB-Funktion replace durch die während der Programmausführung gespeicherten Werte ersetzt. Auf diese Weise werden beispielsweise Prüfdatum, Kunde, Projekt, Prüfer, Prüffrequenz, Spannungsverhältnis und Prüftemperatur automatisiert in das Prüfprotokoll übertragen.

Für die Versuchsergebnisse des Zeit- und Langzeitfestigkeitsbereichs werden die Tabellenzeilen innerhalb von Schleifen automatisch erzeugt. Für jede Probe werden unter anderem Maximalkraft, Minimalkraft, Grundlast, Amplitude und Schwingspielzahl in einer für LaTeX geeigneten Zeichenkette zusammengefasst und anschließend an der entsprechenden Stelle in die Vorlage eingesetzt.

Nach dem Ersetzen aller Platzhalter wird der resultierende LaTeX-Code als neue .tex-Datei gespeichert. Anschließend wird diese über einen Systemaufruf mit pdflatex kompiliert, sodass das Prüfprotokoll ohne weitere manuelle Bearbeitung als PDF-Datei erzeugt wird.